\subsubsection*{\texttt{20-index-token-default}}
\label{app:samples:20_index_token_default}

\paragraph{Q3: How long does the Haste spell last?}
\noindent \textbf{Corr.:}~5.00/5\quad \textbf{Faith.:}~1.00\quad \textbf{Relev.:}~1.00\quad \textbf{Ctx.Rel.:}~1.00\\

\noindent\textbf{Antwort:} The Haste spell lasts for a duration of concentration, up to 1 minute.

\noindent\textbf{Abgerufene Chunks:}
\begin{itemize}[noitemsep, topsep=2pt]
  \item[\textbf{\#1}] \texttt{[0.0167]} \textit{Spells/Haste} -- \#\#\# Haste  *3rd-level transmutation*  **Casting Time:** 1 action  **Range:** 30 \ldots
  \item[\textbf{\#2}] \texttt{[0.0164]} \textit{Spells/Modify Memory} -- \#\#\# Modify Memory  *5th-level enchantment*  **Casting Time:** 1 action  **Range:\ldots
  \item[\textbf{\#3}] \texttt{[0.0161]} \textit{Treasure/Tome of Understanding} -- \#\#\# Tome of Understanding  *Wondrous item, very rare*  This book contains intuit\ldots
  \item[\textbf{\#4}] \texttt{[0.0159]} \textit{Spells/Geas} -- \#\#\# Geas  *5th-level enchantment*  **Casting Time:** 1 minute  **Range:** 60 fee\ldots
  \item[\textbf{\#5}] \texttt{[0.0156]} \textit{Spells/Divination} -- \#\#\# Divination  *4th-level divination (ritual)*  **Casting Time:** 1 action  **R\ldots
\end{itemize}
\vspace{4pt}
\hrule
\vspace{6pt}

\paragraph{Q5: Can a grappled creature cast spells?}
\noindent \textbf{Corr.:}~2.00/5\quad \textbf{Faith.:}~0.00\quad \textbf{Relev.:}~0.00\quad \textbf{Ctx.Rel.:}~—\\

\noindent\textbf{Antwort:} A grappled creature's ability to cast spells isn't addressed in the provided rules for the Arcane Hand spell or the Dash, Disengage, and Dodge actions.

\noindent\textbf{Abgerufene Chunks:}
\begin{itemize}[noitemsep, topsep=2pt]
  \item[\textbf{\#1}] \texttt{[0.0167]} \textit{Monsters/# Monster Statistics} -- the spell slot to do so. For example, a drow mage with the 3rd-level *lightning \ldots
  \item[\textbf{\#2}] \texttt{[0.0164]} \textit{Spells/# Spellcasting} -- \# SPELLCASTING  Magic permeates fantasy gaming worlds and often appears in the f\ldots
  \item[\textbf{\#3}] \texttt{[0.0161]} \textit{Spells/# Spellcasting} -- special abilities that let them cast spells without using spell slots. For examp\ldots
  \item[\textbf{\#4}] \texttt{[0.0159]} \textit{Spells/# Spellcasting} -- images that any creature can see, but the most insidious illusions plant an imag\ldots
  \item[\textbf{\#5}] \texttt{[0.0156]} \textit{Gameplay/Combat} -- enemy.  \#\#\# Two-Weapon Fighting  When you take the Attack action and attack with\ldots
\end{itemize}
\vspace{4pt}
\hrule
\vspace{6pt}

\paragraph{Q12: If a creature has exhaustion level 3 and is also restrained, what is its movement speed?}
\noindent \textbf{Corr.:}~5.00/5\quad \textbf{Faith.:}~1.00\quad \textbf{Relev.:}~1.00\quad \textbf{Ctx.Rel.:}~0.75\\

\noindent\textbf{Antwort:} A creature under the effects of both exhaustion level 3 and the restrained condition has a speed of 0 for the restrained condition. Thus, its movement speed is 0 in this scenario.

\noindent\textbf{Abgerufene Chunks:}
\begin{itemize}[noitemsep, topsep=2pt]
  \item[\textbf{\#1}] \texttt{[0.0167]} \textit{Gamemastering/Conditions} -- suspended, not neutralized.  \#\# Poisoned  - A poisoned creature has disadvantage\ldots
  \item[\textbf{\#2}] \texttt{[0.0164]} \textit{Gamemastering/Conditions} -- | | 2     | Speed halved                                   | | 3     | Disadvant\ldots
  \item[\textbf{\#3}] \texttt{[0.0161]} \textit{Gamemastering/Conditions} -- A frightened creature has disadvantage on ability checks and attack rolls while \ldots
  \item[\textbf{\#4}] \texttt{[0.0159]} \textit{Gameplay/Adventuring} -- creatures moving and the terrain they're moving over.  \#\# Speed  Every character\ldots
  \item[\textbf{\#5}] \texttt{[0.0156]} \textit{Monsters/# Monster Statistics} -- | d20     | 10 1/2             | |              |         |                    |\ldots
\end{itemize}
\vspace{4pt}
\hrule
\vspace{6pt}

\paragraph{Q16: Can you use Action Surge to cast two leveled spells in one turn?}
\noindent \textbf{Corr.:}~2.00/5\quad \textbf{Faith.:}~0.00\quad \textbf{Relev.:}~0.00\quad \textbf{Ctx.Rel.:}~—\\

\noindent\textbf{Antwort:} Based on the new context, the rules for concentrating on spells state that you cannot concentrate on two spells at once. Therefore, even with Action Surge, you cannot cast two leveled spells that require concentration in one turn. However, you could potentially use Action Surge to cast one spell and\ldots

\noindent\textbf{Abgerufene Chunks:}
\begin{itemize}[noitemsep, topsep=2pt]
  \item[\textbf{\#1}] \texttt{[0.0167]} \textit{Classes/Sorcerer} -- | | 5th              | 7                  | |                  |                \ldots
  \item[\textbf{\#2}] \texttt{[0.0164]} \textit{Characterizations/Multiclassing} -- slots, but only to cast your lower-level spells. If a lower-level spell that you\ldots
  \item[\textbf{\#3}] \texttt{[0.0161]} \textit{Classes/Fighter} -- even if the new roll is a 1 or a 2. The weapon must have the two-handed or versa\ldots
  \item[\textbf{\#4}] \texttt{[0.0159]} \textit{Characterizations/Multiclassing} -- from another class.  \#\# Spellcasting  Your capacity for spellcasting depends par\ldots
  \item[\textbf{\#5}] \texttt{[0.0156]} \textit{Classes/Wizard} -- can add your Intelligence modifier to one damage roll of any wizard evocation sp\ldots
\end{itemize}
\vspace{4pt}
\hrule
\vspace{6pt}
